\documentclass[11pt]{article}

\usepackage[margin=1in]{geometry}
\usepackage[super,sort&compress]{natbib}
\usepackage{amsmath}
\usepackage{hyperref}
\usepackage{parskip}
\usepackage{booktabs}
\usepackage{graphicx}

\title{Data-Driven Cardiovascular-Kidney-Metabolic Phenotyping\\
in HCHS/SOL via a Survey-Weighted Dirichlet Process Mixture Model}
\author{}
\date{}

\begin{document}
\maketitle

\section{Rationale}

Hispanic/Latino adults are the largest and fastest-growing minority population in the United States, and cardiovascular disease is their leading cause of death; yet, studied in aggregate, this population appears to have better cardiovascular outcomes than non-Hispanic individuals, a discrepancy that reflects not resilience but the loss of information inherent in treating a heterogeneous, multi-origin population as a single "Hispanic" risk category \citep{blumer2021heterogeneity}. The clearest illustration of that loss is the so-called Hispanic paradox: Hispanic adults have historically shown lower age-adjusted cardiovascular mortality than non-Hispanic White adults despite a substantially worse burden of the risk factors that are supposed to drive that mortality \citep{cortesbergoderi2013paradox,balfour2016cvd}. That paradox dissolves as soon as the aggregate category is disaggregated: cause-specific mortality trends differ sharply by Hispanic/Latino background, with Puerto Ricans showing mortality patterns similar to non-Hispanic Whites while Mexicans show lower rates \citep{rodriguez2017disaggregation}; foreign-born Cubans, Mexicans, and Puerto Ricans have markedly higher cardiovascular mortality than their US-born counterparts \citep{rodriguez2017nativity}; and between 1999 and 2018, stroke mortality rose in Hispanic adults even as it stalled in non-Hispanic White adults, with heart failure mortality accelerating fastest in Hispanic adults under 65 \citep{khan2022trends}. The kidney and metabolic side of the burden is more straightforwardly severe and shows no comparable paradox: within HCHS/SOL itself, chronic kidney disease affects 13.7\% of participants, with more than a two-fold range across Hispanic/Latino background groups (7.4\% in South Americans to 16.6\% in Puerto Ricans) \citep{ricardo2015prevalence}, Hispanic ethnicity is independently associated with a substantially elevated risk of progression to end-stage renal disease even after accounting for diabetes \citep{peralta2006risks,desai2018ckdesrd}, diabetes prevalence in HCHS/SOL ranges from 10.2\% to 18.3\% depending on background \citep{schneiderman2014diabetes}, incident diabetes accrues at 22.1 cases per 1,000 person-years \citep{cordero2022incidence}, and metabolic syndrome prevalence ranges more than two-fold, from 27\% in South Americans to 41\% in Puerto Rican women \citep{heiss2014prevalence}. Aggregating this population, whether into a single "Hispanic" label or, as we argue below, into a single ordinal risk stage, is precisely what has repeatedly been shown to mask the clinically actionable heterogeneity within it \citep{lopezgolden2014disparities}, and unsupervised, data-driven approaches that let subgroups emerge directly from clinical and social data, rather than from investigator-imposed categories, have already begun to reveal cardiometabolic risk profiles that cut across conventional race/ethnicity groupings in the broader US population \citep{ding2023uncovering}. It is against this backdrop, a numerically large, rapidly growing, and demonstrably heterogeneous high-risk population for whom aggregate categories have repeatedly obscured rather than clarified risk, that the adequacy of any single risk-staging tool applied to Hispanic/Latino adults deserves particular scrutiny.

The tool most likely to be applied to this population going forward is Cardiovascular-Kidney-Metabolic (CKM) syndrome, introduced by the American Heart Association (AHA) in 2023 to formalize the pathophysiological interconnection among metabolic risk factors, chronic kidney disease (CKD), and cardiovascular disease (CVD) \citep{ndumele2023synopsis}. Its staging system assigns every individual to one of five ordinal stages, from Stage 0 (no risk factors) to Stage 4 (established clinical CVD). The founding presidential advisory itself frames staging somewhat broadly, as a structure meant to "promote prevention across the life course," explicitly distinct from the companion "prediction algorithms" (the PREVENT equations) that the AHA commissioned separately to furnish continuous cardiovascular risk estimates \citep{ndumele2023advisory}. In practice, however, the staging system has been adopted and validated almost entirely as a risk-stratification device: it is described as designed to "reflect the pathophysiology, spectrum of risk, and opportunities for prevention" \citep{bansal2024framework}, reviewed as a framework that "allows targeted risk stratification and evidence-based management to interrupt CKM trajectories" \citep{gunnarsson2025ckm}, and empirically validated on exactly that basis: in a 14.3-year population-based cohort, CKM stage showed a stepwise, monotonic association with cardiovascular disease and all-cause mortality, leading the authors to conclude that their findings "support the relevance of CKM staging for cardiovascular risk stratification and prevention in the general population" \citep{ahanchi2026longterm}. So whatever the precise balance the AHA advisory struck between staging as a prevention-organizing structure and staging as a graded risk tool, both the field's working interpretation and the subsequent validation literature treat, use, and confirm CKM stage as a risk-stratification instrument, on the premise that stage should track increasing cardiovascular and renal risk and thereby guide stage-specific prevention and treatment. It is precisely this stepwise property between stages that makes the internal composition of any single stage consequential: if a stage is to do stratifying work, the risk within it should be reasonably homogeneous, and a growing body of commentary has argued that CKM staging instead still leaves substantial phenotype-level heterogeneity for stratification to resolve \citep{shen2026editorial}, or that its value lies more in organizing multidisciplinary care than in fine-grained risk discrimination \citep{claudel2023stepstep}. Applied to HCHS/SOL, the population-level pattern that motivates this proposal reappears within the staging tool itself: 50.6\% of the cohort was classified as Stage 2 at baseline \citep{kondeti2026progression}, closely replicating an earlier HCHS/SOL estimate of 53.4\% \citep{chakrabarti2024prevalence}. This is not unique to HCHS/SOL, nor is it a sampling artifact. The same clumping in Stage 2 appears in the Dallas Heart Study (46\%) \citep{shelbaya2025dallas}, in a nationwide Japanese cohort of 266,256 adults (65.0\%) \citep{fujimoto2025stagespecific}, and among postmenopausal women in NHANES (55.7\%) \citep{chen2025postmenopausal}, and a recent review concluded that Stage 2 "represents the most prevalent category, affecting nearly half of adults in Western cohorts" \citep{gunnarsson2025ckm}, a pattern that persists, with variation in degree, across NHANES race/ethnicity strata \citep{ogungbe2024disparities}. The reason is structural rather than incidental: because any CKD, including mild albuminuria with normal estimated glomerular filtration rate, or the presence of a single metabolic risk factor is sufficient to qualify for Stage 2, the criteria function as a broad catch-all \citep{gregg2025riskassessment}, and under the original definition all individuals with type 2 diabetes are uniformly classified as Stage 2 regardless of how well- or poorly-controlled their disease is \citep{gohda2026refining}. A staging system whose purpose is risk stratification loses much of its discriminative value once half of a population, spanning a genuinely wide range of underlying cardiometabolic and renal risk, is compressed into one category, and this is exactly the kind of aggregation that has already been shown, in the broader Hispanic/Latino cardiovascular literature, to mask clinically meaningful subgroup differences rather than reveal them.

This limitation has been explicitly recognized, and several groups have proposed fixes, but each works within the pre-specified AHA staging rather than questioning whether an investigator-imposed ordinal partition is the right representation of cardiometabolic risk to begin with. \citet{tian2026riskbased} developed a sex-specific algorithm to manually subdivide Stage 2 into lower- and higher-risk subgroups (2a/2b), achieving meaningfully different 10-year cardiovascular mortality risk between the two. \citet{gohda2026refining} integrated KDIGO eGFR/albuminuria risk categories into Stage 2 for individuals with type 2 diabetes and showed improved discrimination for CKD progression relative to the unmodified stage. \citet{masoumi2026severity} proposed a continuous CKM severity score as a complement to categorical staging, explicitly to mitigate the limitations of categorical staging, and demonstrated a substantial mortality-risk gradient within Stage 2 alone. Two unsupervised approaches have begun to probe this heterogeneity empirically: \citet{yang2025datadriven} applied $k$-means clustering with a fixed number of clusters to UK Biobank participants in CKM Stages 1--3 and recovered five phenogroups with materially different long-term risk of stroke, chronic kidney disease, and liver disease, while \citet{zhou2025metabolomic} used untargeted plasma metabolomics and unsupervised clustering to identify metabolic subtypes cross-cutting AHA-defined stages. A parallel commentary literature has begun to question the added value of the CKM construct itself relative to existing risk tools, asking whether it represents "a new frontier or simple rebranding" \citep{mannegoehler2025rebranding} and flagging specific structural gaps such as the omission of hepatic biomarkers despite their established role upstream of insulin resistance \citep{cai2025hype}. Every one of these corrective efforts, however, is supervised by the AHA staging system: each starts from the pre-defined stages, most often Stage 2 alone, and attempts to split or score risk within that boundary, rather than asking what structure the joint distribution of cardiovascular, kidney, and metabolic biomarkers would reveal if no ordinal staging rule were imposed at all. Where unsupervised methods have been used, the number of clusters has been fixed a priori and the cohorts studied, UK Biobank, single-center metabolomics samples, do not reflect the ethnic and cardiometabolic heterogeneity documented within the U.S. Hispanic/Latino population and summarized above. A prior latent class analysis in HCHS/SOL identified four subclasses of metabolic syndrome expression \citep{arguelles2015characterization}, demonstrating that this cohort's biomarker data do contain recoverable substructure, but that analysis used a small, pre-specified set of canonical metabolic syndrome components rather than the fuller panel of cardiac, renal, and metabolic biomarkers now organized under the CKM framework, used a parametric latent class model that still required the number of classes to be fixed and compared across candidate values, and was not linked to, or evaluated against, the subsequently introduced AHA CKM staging.

This motivates the present proposal: to derive cardiovascular-kidney-metabolic phenotypes for the entire HCHS/SOL cohort (n\,=\,16,415) directly from the joint distribution of CKM-relevant biomarkers, with no ordinal staging structure imposed and no number of phenotypes fixed in advance, and only afterward to compare the resulting data-driven phenotype assignments against each participant's AHA CKM stage. If Stage 2 truly represents a biologically coherent, homogeneous risk category, the data-driven phenotypes should recover that homogeneity. If, as the sub-staging and continuous-score literature above suggests, Stage 2 instead pools together individuals with materially different multi-organ risk profiles, an unsupervised approach applied to the whole cohort should reveal that structure directly, without requiring an analyst to first guess which covariates might split the stage, and should do so in the one large, richly phenotyped cohort constructed specifically to represent the diversity of the U.S. Hispanic/Latino population rather than a single background group or a non-Hispanic reference population.

Two classes of methods dominate the existing cardiometabolic phenotyping literature: hard-partition algorithms such as $k$-means, used by \citet{yang2025datadriven}, and finite mixture or latent class models with a pre-specified or model-selected number of classes, used throughout the metabolic syndrome latent class analysis literature, including within HCHS/SOL itself \citep{arguelles2015characterization,riahi2018patterns}. Both require the number of clusters $K$ to be fixed, or chosen from a discrete candidate set via information criteria, before the clustering step is run. A Dirichlet process mixture model (DPMM) relaxes this requirement: it places a nonparametric prior over the partition of participants into clusters, allowing the number of supported phenotypes to be estimated jointly with the phenotypes themselves, and it returns full posterior uncertainty over cluster assignment rather than a single hard partition \citep{dahl2006modelbased}. DPMMs have an established track record for exactly this kind of task, unsupervised phenotype discovery from a multivariate clinical marker profile without a pre-specified number of subtypes, including symptom-based phenotyping in Parkinson's disease \citep{white2012dirichlet}, profile-regression-style clustering of clinical covariates \citep{liverani2013premium}, and, most directly relevant, clustering of diabetes multimorbidity from electronic health record data, where DPMM was shown to adapt its cluster count to underlying data structure and to outperform fixed-$K$ alternatives on high-dimensional, imbalanced clinical marker sets \citep{kita2025enhancing}. A concurrent, independently developed data-driven phenogroup classification of cardiometabolic risk \citep{feng2026phenogroups} further indicates that expanding beyond the canonical metabolic syndrome components to a richer risk-factor panel reveals phenogroups with materially different cardiovascular hazard, reinforcing that the CKM biomarker panel is rich enough to support this kind of discovery.

HCHS/SOL was not collected as a simple random sample: it is a two-stage, stratified, clustered area-probability sample of Hispanic/Latino households across four U.S. field centers, with unequal probabilities of selection by design, requiring dedicated sampling-weighted methodology even for standard analyses conducted within the cohort \citep{sorlie2010design,lavange2010sample}. \citet{lin2014genetic} demonstrate explicitly, in this same cohort, that ignoring the complex sampling design in model-based analyses of HCHS/SOL data can bias inference, motivating design-aware estimation even for genetic association models. A DPMM fit naively to the observed HCHS/SOL sample, ignoring the sampling weights, would let over- or under-sampled strata, by center, background group, and other design strata, distort both the estimated number of phenotypes and the estimated size and biomarker profile of each phenotype, the same design-induced bias that has been documented broadly for Bayesian nonparametric models fit to complex survey samples without weighting \citep{si2013bayesian,kunihama2014nonparametric}. We therefore adopt a pseudolikelihood formulation in which each participant's contribution to the DPMM likelihood is exponentiated by their normalized survey weight, so that the fitted mixture targets the sampling-weighted Hispanic/Latino population represented by HCHS/SOL rather than merely the realized sample \citep{savitskytoth2015bayesian,williamssavitsky2018unattenuated}. Exponentiating the likelihood by weights alone, however, is insufficient for valid inference: the within-cluster dependence induced by a multistage design inflates the effective information contributed by each unit relative to what plain weighting assumes, which in turn produces under-coverage of pseudo-posterior credible regions if left uncorrected. We address this using the sandwich-type Godambe information correction of \citet{williamssavitsky2021uncertainty}, which applies a fast, one-time post-processing projection to Markov chain Monte Carlo draws from the pseudo-posterior, rescaling and reshaping the pseudo-posterior covariance matrix so that nominal coverage of cluster-parameter uncertainty is approximately restored. This correction is implemented in the \texttt{csSampling} framework \citep{hornby2023cssampling} and has since been extended to multilevel model structures \citep{williamsmcguiresavitsky2025uncertainty}, making it directly applicable to a hierarchically clustered design such as HCHS/SOL's. In combination, the survey-weighted pseudolikelihood and Godambe-adjusted pseudo-posterior allow us to recover phenotypes, phenotype prevalences, and posterior uncertainty on cluster membership that are simultaneously representative of the target Hispanic/Latino population rather than the realized, design-biased sample, and accompanied by correctly calibrated uncertainty, rather than the artificially narrow credible intervals that a naive weighted pseudo-posterior would otherwise produce.

Applying an unweighted, fixed-$K$ clustering algorithm, or simply re-deriving sub-strata within the existing AHA Stage 2 definition, would leave open the question motivating this proposal: is the concentration of half of the HCHS/SOL cohort in a single CKM stage a real reflection of homogeneous risk in this population, or an artifact of an ordinal staging rule that was not designed with the joint multivariate structure of cardiovascular, kidney, and metabolic biomarkers, or with a heterogeneous, multi-origin, historically underrepresented population, in mind? By phenotyping the entire cohort from the biomarker data directly, with the number of phenotypes determined by the data rather than fixed by the analyst, and with survey design and uncertainty properly accounted for, and only subsequently cross-tabulating the resulting phenotypes against AHA CKM stage, this proposal provides a direct, population-representative empirical test of that question. Given the documented within-Stage-2 risk heterogeneity already reported using stage-conditioned sub-classification and continuous scoring approaches \citep{tian2026riskbased,masoumi2026severity}, and the burden and heterogeneity in cardiovascular, kidney, and metabolic risk across Hispanic/Latino background groups documented above \citep{heiss2014prevalence,ricardo2015prevalence,lopezgolden2014disparities}, we anticipate that a fully data-driven phenotyping of HCHS/SOL will recover clinically distinct subgroups currently pooled under Stage 2. Because this population already bears a disproportionate and heterogeneous cardiovascular-kidney-metabolic disease burden that aggregate categories have repeatedly been shown to mask, and because it remains underrepresented in the cohorts that originally informed the AHA staging framework, resolving whether that framework's coarsest stage reflects real homogeneity or an artifact of its design has direct and immediate implications for precision cardiovascular-kidney-metabolic prevention in Hispanic/Latino adults specifically.

\paragraph{Optional analysis: incorporating metabolomics.} Three batches of metabolomics data exist within HCHS/SOL, but, having been generated as part of an ancillary study rather than the core examination, they are available for only a subset of participants rather than the full cohort. This partial coverage should not be allowed to constrain the primary analysis: the survey-weighted DPMM described above is fit on the clinical cardiovascular, kidney, and metabolic biomarker panel that is available for essentially the entire cohort, precisely so that every participant receives a phenotype assignment regardless of ancillary-study enrollment. Metabolomics is instead treated as an optional, secondary layer, for two reasons beyond incompleteness. First, three separate measurement batches introduce technical variation that, left uncorrected, can itself masquerade as biological clustering structure; before any downstream use, batch-level correction is required, using quality-control-sample-based signal correction of the kind developed for large-scale, multi-batch metabolic profiling studies \citep{dunn2011procedures}, a removal-of-unwanted-variation approach designed specifically for combining metabolomics data across batches and platforms \citep{delivera2012normalizing}, or a comparable benchmarked batch-correction method \citep{wehrens2016improved}; batch should also be retained as a covariate and checked against the final phenotype assignments. Second, ancillary-study enrollment is itself a selection mechanism that need not be random with respect to CKM phenotype, so any downstream use of the metabolomics subsample should not simply assume it is missing completely at random.

Two complementary strategies allow metabolomics to be incorporated once these prerequisites are addressed, and both preserve the primary, full-cohort phenotyping rather than replacing it. The simpler strategy is post hoc validation: within the metabolomics-available subsample, cross-tabulate the phenotype assignments already obtained from the biomarker-only DPMM against metabolomic profiles to test whether any recovered phenotype, particularly one overlapping the coarse AHA Stage 2 category, is further separable on metabolomic grounds, mirroring the logic already used by the CKM sub-staging literature discussed above but applied to data-driven rather than investigator-defined phenotypes. The more ambitious strategy is to extend the joint model itself: rather than restricting analysis to complete cases or imputing metabolomics for the majority of the cohort that never had it measured, a further block can be added to the DPMM likelihood whose contribution is included only for participants with observed metabolomics and simply omitted, not imputed, for everyone else, following the general approach to Bayesian integrative modeling of multi-omics data with block-wise missingness \citep{fang2018bayesian,lin2020general}. A related and more explicitly nonparametric alternative is to let the metabolomics subsample support its own clustering that adheres loosely to, rather than being forced to coincide with, the consensus phenotyping obtained from the full-cohort biomarker panel, in the spirit of Bayesian consensus clustering for multiple, unequally covered data sources \citep{lock2013bayesian} and its recent extension to a fully nonparametric, two-level Dirichlet process construction that similarly avoids prespecifying the number of clusters in either source \citep{burhanuddin2026integrative}; partial multi-omics clustering methods developed for the analogous problem of cancer subtyping with incomplete omics coverage, which explicitly avoid imputing missing platforms, offer a further precedent \citep{rappoport2018nemo}. Naively imputing metabolomics values for the full cohort before clustering is deliberately avoided under either strategy, since imputation-based approaches to missing rows in multi-omics integration have been shown to produce rapidly growing uncertainty in the resulting structure as the proportion of missingness increases \citep{voillet2016handling}, which is precisely the regime an ancillary substudy occupies. Finally, because ancillary-study membership constitutes a second, non-random sampling phase nested within the primary HCHS/SOL design, any analysis of the metabolomics subsample should reweight participants by the inverse probability of ancillary-study selection, following established two-phase survey methodology \citep{breslow2009improved} and approaches for borrowing information from auxiliary data available on only part of a cohort \citep{chen2021improving}, combined multiplicatively with the primary HCHS/SOL design weight rather than substituted for it; comparable inverse-probability-weighting corrections have been shown to materially reduce selection bias in omics analyses restricted to non-randomly available subsamples \citep{carry2021inverse}.

\section{Methods: Simulation Validation Study}

Before the survey-weighted DPMM is fit to HCHS/SOL itself, we validate the pseudolikelihood-plus-Godambe-correction approach in a fully controlled simulation where the truth is known, extending the informative-sampling designs of \citet{williamssavitsky2021uncertainty} (their PPS1 and PPS3 mechanisms) from a single continuous outcome to a $K_0$-component multivariate Gaussian mixture outcome standing in for a CKM-marker-like biomarker panel. This subsection describes the original validation design as implemented in \texttt{simulations/archive/src/R/dpmm\_survey\_simulation.R}; the current, actively developed pipeline (data-generating process, fitting, and correction) lives in \texttt{simulations/src/R/} and is described in the subsequent subsections, which supersede several choices made here.

\paragraph{Finite population and informative design.} We generate a finite population of $N$ elements nested in $n_{\text{psu}}$ primary sampling units (PSUs), which are in turn nested in $H$ fixed strata under proportional allocation. Each element carries an auxiliary variable $x_2 > 0$ that plays two roles simultaneously, reproducing the PPS1 informative-selection mechanism: it drives true mixture-component membership through a multinomial-logit model, $\Pr(z_i = k \mid x_{2i}) \propto \exp(\alpha_k + \beta_k x_{2i})$, and it drives the unit's probability of selection into the sample, $\pi_i \propto x_{2i}$, so that composition and sampling probability are confounded by design. Each PSU additionally carries an independent effect $z_2$ (an exponential deviate, centered), which is added to the outcome only and never enters the selection mechanism, reproducing the PPS3 mechanism by which cluster-level dependence inflates effective variance without itself being informative. Conditional on true component membership $z_i = k$, the $P$-dimensional outcome is drawn as $Y_i \sim \mathcal{N}_P(\mu^0_k, \Sigma^0_k) + z_{2,\text{psu}(i)}$ for a fixed set of $K_0$ superpopulation truth parameters $(\mu^0_k, \Sigma^0_k)$.

\paragraph{Sampling.} A stratified, approximately-PPS sample of target size $n$ is drawn by allocating $n$ proportionally to stratum population size and then applying independent Poisson-PPS selection within each stratum on that stratum's own $x_2$-based inclusion probabilities, $\pi_i \propto x_{2i}$ (capped below 1). Each sampled unit's design weight is $w_i = 1/\pi_i$, normalized to sum to the realized sample size. Stratum and PSU membership are retained for every sampled unit, since both are required later for the design-based variance correction.

\paragraph{Survey-weighted truncated mixture.} The DPMM is fit via a blocked Gibbs sampler on a truncated representation with truncation level $L$. Component parameters $(\mu_k, \Sigma_k)$ are updated via a weighted conjugate Normal-Inverse-Wishart step in which each unit's contribution to the sufficient statistics is scaled by its normalized survey weight $w_i$, with a minimum-eigenvalue floor imposed on each $\Sigma_k$ for numerical stability, and the mixture weights $\pi_k$ are likewise updated from weighted component-occupancy counts under a sparse Dirichlet$(\alpha_0/L,\dots,\alpha_0/L)$ prior with $\alpha_0$ small, the overfitted-mixture construction of \citet{rousseau2011asymptotic}. Component indicators, by contrast, are updated \emph{without} the per-unit weight exponent, $\Pr(z_i = k) \propto \pi_k f(Y_i \mid \mu_k, \Sigma_k)$, rather than the fully weighted pseudolikelihood form $\Pr(z_i = k) \propto [\pi_k f(Y_i \mid \mu_k, \Sigma_k)]^{w_i}$ used in the closely related survey-weighted overfitted latent class literature \citep{stephenson2024wolca,wu2024swolca}. This departure is the result of a diagnostic investigation, not an a priori choice: exponentiating the individual-unit classification likelihood by $w_i$ was found to induce a classification-instability pathology under this project's continuous, covariate-driven (PPS) informative-sampling design, in which units with below-median weight are reassigned to a different component in roughly 40\% of Gibbs iterations regardless of chain length, with the truncated components' occupancy never concentrating on the true $K_0$ and instead remaining spread, with little decay, across all $L$ truncated slots. This behavior persisted under a quasi-Bernoulli stick-breaking overlay \citep{zeng2020consistent}, under the sparse Dirichlet prior above on its own, under \citeauthor{stephenson2024wolca}'s and \citeauthor{wu2024swolca}'s own published MCMC scale ($L=20$--$30$, 10,000--20,000 iterations), and even when the population/sample was regenerated to match their own validated survey design exactly (two strata with stratum-dependent class membership and a constant weight per stratum, rather than this project's continuous $x_2$-driven design); it was resolved only by removing the per-unit weight exponent from the classification step while retaining $w_i$ in every aggregate update. The two survey-weighted overfitted latent class papers above use $w_i$-exponentiated classification for a categorical-manifest-variable, supervised-outcome setting; an earlier method from the same research group, applied to dietary pattern discovery in HCHS/SOL itself, instead fits its clustering model fully unweighted and applies the survey design only post hoc, to reweight the fitted classification into a population-representative estimate \citep{stephenson2020rpchchssol}, an even more design-agnostic choice at estimation time than the one adopted here. Two fitted-model arms are compared, mirroring the design-aware/design-oblivious contrast used to validate PPS1 in \citet{williamssavitsky2021uncertainty}: an \emph{informed} arm, in which the mixture weights are allowed to depend on $x_2$ through a covariate-dependent construction (each component's weight fit as a weighted logistic regression on $x_2$, shrunk toward the sparse-Dirichlet baseline, a practical stand-in for full P\'olya-Gamma augmentation), and a \emph{naive} arm, in which a single global sparse Dirichlet prior is used and the design variable is omitted from the fitted model entirely, reproducing the design-ignoring analyst who fits an unweighted or covariate-blind mixture to informatively sampled data. This project's \texttt{baysc} companion diagnostics \citep{wu2026baysc} confine themselves to categorical manifest variables and were not usable directly for this continuous, multivariate-Gaussian CKM-marker setting; the comparison above instead re-implements their published weighting mechanism on this project's own outcome model to isolate the weighting choice as the variable under test.

\paragraph{Label-switching-safe prevalence functional.} Because mixture component labels are not identified across MCMC iterations, we do not attempt to track $(\mu_k, \Sigma_k)$ directly across draws. Instead, post-burn-in draws of the hard cluster assignment are collapsed to the $K_{\text{active}}$ most frequently occupied components and relabeled using the Stephens (2000) algorithm \citep{stephens2000dealing}, as implemented in the \texttt{label.switching} package, applied to a one-hot encoding of the collapsed assignments. Because $L > K_0$, $K_{\text{active}}$ is generally larger than $K_0$ even after the classification-step correction above; the $K_0$ fitted components with centroids nearest the known true phenotype means are matched one to one to the $K_0$ phenotypes, and every remaining fitted component is merged into whichever true phenotype it is nearest to, rather than discarded, so that the resulting prevalence vector sums to one. This merge-to-known-truth step is used for the validation results in this subsection, where $K_0$ and the true component means are known by construction and can therefore be used to score the method; on real HCHS/SOL data, with $K_0$ unknown, this step is replaced by an unsupervised, data-driven merge, described and validated later in this section. For each relabeled and merged draw, the survey-weighted phenotype prevalence vector $\hat\pi = (\hat\pi_1, \dots, \hat\pi_{K_0})$ is computed as the weight share falling in each merged phenotype, $\hat\pi_k = \sum_i w_i \mathbf{1}(z_i = k) / \sum_i w_i$. This relabeled, merged prevalence vector, rather than the raw component parameters, is the estimand carried into the uncertainty correction below.

\paragraph{Godambe sandwich correction on the prevalence functional.} Exponentiating the likelihood by design weights alone understates posterior uncertainty because it ignores the within-PSU dependence and unequal selection probabilities induced by the multistage design. We apply the sandwich-type adjustment of \citet{williamssavitsky2021uncertainty} (their Algorithm~1) to the permutation-invariant phenotype prevalence functional. While an initial hard-classification estimating equation conditioned on a single modal assignment vector and understated classification-boundary ambiguity, the validated pipeline constructs the Godambe pieces from the \emph{marginal mixture pseudolikelihood}:
\begin{equation}
m_i = \sum_{k=1}^{K_0} \hat\pi_k a_{ik}, \quad u_{ik} = w_i \frac{a_{ik} - a_{i, K_0}}{m_i}, \quad \hat H = \sum_{i=1}^n w_i \frac{(a_i - a_{i, K_0})(a_i - a_{i, K_0})^\top}{m_i^2}
\end{equation}
where $a_{ik}$ is the sub-mixture density for phenotype $k$ reconstructed from the model's unmerged raw fitted components mapped to phenotype $k$. The design-based score covariance $\hat J$ is obtained from a stratified, PSU-level with-replacement bootstrap of the weighted score total \citep{raowu1988resampling}, resampling PSUs independently within each stratum across $B = 200$ replicates, giving the marginal-likelihood sandwich piece $\hat\Sigma_{\text{sandwich}} = \hat H^{-1}\hat J \hat H^{-1}$. This sandwich piece alone, however, only reflects design/classification variability given the fitted component parameters $(\mu_k, \Sigma_k)$ treated as fixed; it omits the additional variability contributed by $(\mu_k,\Sigma_k)$-mixing uncertainty within the MCMC chain itself. We therefore combine $\hat\Sigma_{\text{sandwich}}$ with the raw pseudo-posterior covariance $\hat\Sigma_{\text{naive}} = \text{Cov}(\hat\pi_m)$ additively, treating the two as (approximately independent) components of total variance:
\begin{equation}
\hat\Sigma_{\text{corrected}} = \hat\Sigma_{\text{sandwich}} + \hat\Sigma_{\text{naive}}
\end{equation}
and impose $\hat\Sigma_{\text{corrected}}$ on the posterior draws of $\hat\pi$ via Cholesky rotation:
\begin{equation}
\hat\pi_m^a = (\hat\pi_m - \bar\pi) R_2^{-1} R_1 + \bar\pi
\end{equation}
where $R_1^\top R_1 = \hat\Sigma_{\text{corrected}}$ and $R_2^\top R_2 = \hat\Sigma_{\text{naive}}$. In a pilot comparison ($n=200$ replicates), using $\hat\Sigma_{\text{sandwich}}$ alone in place of $\hat\Sigma_{\text{corrected}}$ -- the marginal-likelihood correction without this total-variance recombination step -- left the nominal-interval coverage of every phenotype below that obtained with the additive correction (74--84\% versus 85--89\%), motivating its inclusion in the validated pipeline reported in Table~\ref{tab:swolca-ours-1000}.

\paragraph{Head-to-head comparison against the survey-weighted overfitted LCA classification rule.} To confirm that the classification-instability finding above reflects the per-unit weighting mechanism itself, rather than an artifact of this project's harder continuous PPS design, we re-ran the comparison entirely on \citeauthor{wu2024swolca}'s own validated population/sample regime: two strata with class-membership probabilities that differ by stratum (the sole informativeness mechanism) and a constant design weight per stratum, $N_s/n_s$, with no continuous selection-probability covariate at all, following their published data-generating code. Both weighting rules, $w_i$-exponentiated classification (their published mechanism, applied here to this project's continuous-outcome mixture rather than their categorical manifest-variable model) and unweighted classification with weighted aggregate updates (this project's fix), were fit to identical population draws, with everything else, including the sparse Dirichlet mixture-weight prior, held fixed. Over 30 replicates per method (Table~\ref{tab:swolca-compare}), the $w_i$-exponentiated classification rule reproduced the same fragmentation pathology even on its own native design (mean $K_{\text{active}}$ recovered at truncation, $10$ of $L=10$, versus $6.0$ under the unweighted rule) and showed substantially larger relative bias on two of the four phenotype prevalences (up to $16.8\%$, versus at most $1.7\%$ under the unweighted rule). We take this as evidence that the instability is attributable to the individual-level weighting mechanism itself under this class of continuous-outcome truncated mixture models, not to the specific informative-sampling design used to validate it, and it is consistent with the fact that an earlier method from the same research group, applied to HCHS/SOL directly, avoided per-unit weighting during estimation entirely \citep{stephenson2020rpchchssol}.

\begin{table}[h]
\centering
\small
\resizebox{\linewidth}{!}{%
\begin{tabular}{lrrrr}
\toprule
& Phenotype 1 & Phenotype 2 & Phenotype 3 & Phenotype 4 \\
\midrule
True prevalence & 0.1875 & 0.1625 & 0.2875 & 0.3625 \\
\midrule
$w_i$-exponentiated classification: relative bias & $+16.2\%$ & $-16.8\%$ & $-1.2\%$ & $+0.1\%$ \\
Unweighted classification (ours): relative bias & $-1.7\%$ & $+0.9\%$ & $-0.4\%$ & $+0.8\%$ \\
\bottomrule
\end{tabular}%
}
\caption{Relative bias in estimated phenotype prevalence, $n=30$ replicates per method, on \citeauthor{wu2024swolca}'s own two-strata population/sample regime. Mean recovered active-component count: $10.0$ of $L=10$ ($w_i$-exponentiated) versus $6.0$ (unweighted).}
\label{tab:swolca-compare}
\end{table}

Scaling the unweighted-classification (ours) arm to $n = 1{,}000$ independent Monte Carlo replicates under the raw-mixture marginal Godambe adjustment with total-variance recombination (Table~\ref{tab:swolca-ours-1000}) confirms that prevalence estimation is essentially unbiased and substantially, though not perfectly, closes the uncertainty-calibration gap found in every earlier formulation. Across all four phenotypes, point estimation bias is negligible ($|\text{Bias}| \le 0.007$, with Phenotype~3 showing zero bias to four decimal places). Estimated standard errors implied by the corrected sandwich track the empirical Monte Carlo standard errors far more closely than under the hard-classification or marginal-likelihood-only formulations, with SE ratios ($\text{Empirical SE} / \text{Estimated SE}$) ranging from $1.16$ to $1.49$ (versus $1.47$--$2.20$ for the hard-classification rule). Standard coverage of the nominal 95\% Godambe-adjusted interval ranges from $84.7\%$ to $91.8\%$; decomposing this coverage by centering intervals at the empirical replicate mean $\bar\pi$ (Test~A, which isolates the pure variance calibration from residual location error) raises this to $93.0\%$--$94.1\%$ across all four phenotypes. At $n=1{,}000$ replicates the Monte Carlo standard error of a coverage proportion near the nominal rate is $\sqrt{0.95 \times 0.05 / 1000} \approx 0.69$ percentage points, so these Test~A values remain a small but statistically detectable $1.3$--$2.9$ standard errors below nominal, not fully within simulation noise. We report this honestly as a substantial reduction of, rather than a complete resolution to, the coverage gap; the most likely remaining source, not yet corrected for, is posterior uncertainty in the fitted component parameters $(\mu_k, \Sigma_k)$ themselves, which the plug-in construction of $a_{ik}$ above does not propagate.

\begin{table}[h]
\centering
\small
\resizebox{\linewidth}{!}{%
\begin{tabular}{lrrrrrrr}
\toprule
Phenotype & True $\pi_0$ & Bias & Estimated SE & Empirical SE & SE Ratio & Standard 95\% Cov. & Centered Cov. (Test A) \\
\midrule
1 & 0.2747 & $-0.0022$ & 0.0173 & 0.0225 & 1.30 & 87.9\% & 93.0\% \\
2 & 0.2600 & $-0.0047$ & 0.0138 & 0.0207 & 1.49 & 84.7\% & 93.2\% \\
3 & 0.2394 & $\hphantom{-}0.0000$ & 0.0144 & 0.0166 & 1.16 & 91.8\% & 93.9\% \\
4 & 0.2260 & $+0.0070$ & 0.0173 & 0.0216 & 1.25 & 89.8\% & 94.1\% \\
\bottomrule
\end{tabular}%
}
\caption{Validation results for the survey-weighted DPMM with raw-mixture marginal Godambe adjustment ($n=1{,}000$ Monte Carlo replicates). SE Ratio is Empirical~SE / Estimated~SE; Centered Coverage evaluates the nominal 95\% interval centered at the empirical mean $\bar\pi$ to isolate variance calibration from residual finite-sample location error.}
\label{tab:swolca-ours-1000}
\end{table}

\paragraph{Validation metrics and execution.} Each replicate proceeds end to end: generate population, draw the PPS sample, fit the weighted truncated mixture, relabel/merge and summarize the prevalence functional, and apply the Godambe correction. Validation evaluates the fitted model against the known superpopulation truth on empirical coverage of the 95\% Godambe-adjusted credible interval, relative bias of $\hat\pi$, raw active components $K_{\text{raw}}$, and numerical stability of the covariance updates. Replicates are executed as independent SLURM array tasks with deterministic seeds. Under the corrected unweighted-classification rule with the sparse Dirichlet mixture-weight prior and the raw-mixture marginal sandwich adjustment with total-variance recombination, the methodology demonstrates essentially unbiased finite-population recovery and substantially, though not completely, corrected inferential calibration under complex informative sampling, with a small residual undercoverage documented above as an open item for further refinement.

\paragraph{Revised data-generating process: clustering in phenotype-membership probabilities.} The single-arm result above (Table~\ref{tab:swolca-ours-1000}) was validated on a design in which cluster/stratum heterogeneity is injected as an additive shift directly on the biomarker outcome, $Y_i = \mu^0_{z_i} + c(a_h-1) + c(b_{hi}-1) + \epsilon_i$. Extending this design to a two-point clustering-strength grid ($c \in \{0.3, 1\}$) surfaced a structural problem with that construction: because the cluster shift enters $Y$ directly, units in strongly-shifted clusters are pulled toward a different phenotype's biomarker profile than their true latent class, producing cluster-correlated misclassification that no post hoc variance correction can remove. This showed up empirically as the raw number of active mixture components, $K_{\text{raw}}$, growing with clustering strength (from $6.3$ at $c=0.3$ to $8.5$ at $c=1$) rather than remaining stable, since fitted components were being spent to absorb cluster identity rather than representing only true phenotype structure.

We therefore relocated the clustering mechanism from the outcome to the phenotype-membership probabilities, in keeping with the general principle, already present in the survey-weighted latent-class literature, that design effects act on class-membership probabilities rather than on the manifest variables' conditional distribution \citep{stephenson2024wolca,wu2024swolca}. Concretely, $Y_i \mid z_i = k \sim \mathcal N_P(\mu^0_k,\Sigma^0_k)$ is kept invariant across clusters, and stratum/cluster heterogeneity instead enters the multinomial-logit membership probability,
\begin{equation}
\Pr(z_i = k \mid x_{2i}, h, \text{psu}(i)) \propto \exp\!\big[(x_{2i}-2)\beta_k + c\,(a_{hk} + b_{hi,k})\big],
\end{equation}
where $a_{hk}$ and $b_{hi,k}$ are stratum- and cluster-level random effects on the $k$-th class logit, each centered to sum to zero across phenotypes within its stratum/cluster so that $c$ governs clustering strength without shifting marginal phenotype prevalences. Under this revised design (implemented in \texttt{simulations/src/R/generate\_population.R}), validated at $c=0.3$, $K_{\text{raw}}$ is stable at $6.0$, confirming that the fitted mixture components once again represent only true phenotype structure.

\paragraph{Two further refinements to the Godambe correction.} First, because $\hat\Sigma_{\text{sandwich}}$ is a cluster-robust variance estimator, it is known to be biased downward when the number of sampled PSUs is small relative to the number of strata. We correct this with a finite-sample design-degrees-of-freedom adjustment: $\hat\Sigma_{\text{sandwich}}$ is scaled by $c_{\text{corr}} = M/(M-H)$ ($M$ sampled PSUs, $H$ strata) before being combined with $\hat\Sigma_{\text{naive}}$, and the credible interval is built from a $t_{M-H,\,0.975}$ critical value in place of $z_{0.975}$. Second, the design-based score covariance $\hat J$, originally estimated via a $B=200$-replicate stratified PSU bootstrap, is now obtained from the closed-form, delete-one-PSU jackknife: for the linear (additive-over-PSU) score total used here, the delete-one-PSU jackknife variance is algebraically identical to the classical Rao-Wu ``ultimate cluster'' with-replacement variance estimator \citep{raowu1988resampling}, so this replaces a Monte Carlo approximation with its own exact, deterministic limit, at lower computational cost and with no resampling noise.

\paragraph{A third refinement: eliminating a selection bias in the raw pseudo-posterior covariance.} The relabeling step (above) collapses each post-burnin MCMC iteration's raw component labels onto the $K_{\text{active}}$ most-occupied components and, in its original form, discarded the \emph{entire} iteration if even one of the $\approx 2{,}000$ sampled units carried a label outside that active set. Across replicates this discarded, on average, more than half of the $1{,}000$ retained post-burnin iterations (as few as $16\%$ survived in some replicates), and an autocorrelation-adjusted effective-sample-size check (\texttt{coda::effectiveSize}) on the surviving draws showed the loss compounded further, to as little as $5$--$15\%$ of the nominal $1{,}000$ draws for some phenotypes. Critically, the discarded iterations were not a random subsample: they were systematically the ones in which at least one unit's classification was atypical, so $\hat\Sigma_{\text{naive}} = \text{Cov}(\hat\pi_m)$, computed only from the surviving iterations, was subject to a selection bias rather than merely reduced precision. We corrected this by remapping each out-of-set label to its nearest active component, by Euclidean distance between the final-iteration fitted centroids, rather than discarding the iteration entirely, so every one of the $1{,}000$ retained post-burnin iterations now contributes to $\hat\Sigma_{\text{naive}}$.

Tables~\ref{tab:cluster-probs-uncorrected} and~\ref{tab:cluster-probs-corrected} report validation metrics under the revised phenotype-probability clustering design at $c=0.3$, with all three refinements (jackknife, design-df correction, full-chain retention) except where noted, without and with the design-degrees-of-freedom correction ($n=1{,}000$ replicates). Relative to the additive-shift design (Table~\ref{tab:swolca-ours-1000}), coverage improves substantially even before the design-df correction is applied, and that correction adds a further, consistent $2$--$3$ percentage points on top: standard 95\% coverage reaches $91.8$--$96.9\%$ across all four phenotypes, with three of four phenotypes at or above the 95\% nominal target once the design-df correction is applied. Point estimation remains essentially unbiased ($|\text{bias}| \le 0.0044$).

\paragraph{Testing, and ruling out, two candidate explanations for the residual gap.} A modest coverage shortfall remains on some phenotypes even with all three refinements. We tested two specific, falsifiable explanations for it, rather than asserting a cause without evidence. First, we tested whether the gap is attributable to unpropagated posterior uncertainty in the fitted mixture components $(\mu_k,\Sigma_k)$, by rebuilding the marginal-likelihood density $a_{ik}$ from the \emph{true} superpopulation Gaussian densities instead of the fitted ones, holding everything else in the pipeline fixed (an oracle substitution, re-verified after the full-chain-retention fix above). Coverage under this oracle substitution was no better, and on some phenotypes slightly worse, than under the fitted components ($89.9$--$93.1\%$ versus $91.8$--$95.6\%$), so this explanation is not supported. Second, we tested whether the residual gap reflects Monte Carlo imprecision in $\hat\Sigma_{\text{naive}}$ by doubling the MCMC chain length ($3{,}000$ iterations, $1{,}000$ burn-in, versus $1{,}500$/$500$); this raised effective sample size by $30$--$90\%$ per phenotype but produced no corresponding coverage improvement, and one phenotype's corrected-arm coverage fell by several points, within the range expected from re-running an independent $1{,}000$-replicate Monte Carlo experiment. Both tests point the same direction: the residual gap is not primarily a precision or plug-in-uncertainty problem in the variance construction, and its specific source remains an open question for further work.

\begin{table}[h]
\centering
\small
\resizebox{\linewidth}{!}{%
\begin{tabular}{rrrrrr}
\toprule
Phenotype & True $\pi_0$ & Bias & SE Ratio & Standard 95\% Cov. & Centered Cov. (Test A) \\
\midrule
1 & 0.2706 & $-0.0006$ & 1.06 & 91.8\% & 91.7\% \\
2 & 0.2579 & $-0.0044$ & 1.06 & 91.8\% & 93.6\% \\
3 & 0.2418 & $\hphantom{-}0.0014$ & 0.94 & 93.9\% & 93.5\% \\
4 & 0.2297 & $+0.0036$ & 1.18 & 95.6\% & 96.1\% \\
\bottomrule
\end{tabular}%
}
\caption{Validation results under the revised phenotype-probability clustering design ($c=0.3$), jackknife-based Godambe correction with full-chain retention, \emph{without} the design-degrees-of-freedom adjustment ($n=1{,}000$ Monte Carlo replicates). Mean $K_{\text{raw}} = 6.0$.}
\label{tab:cluster-probs-uncorrected}
\end{table}

\begin{table}[h]
\centering
\small
\resizebox{\linewidth}{!}{%
\begin{tabular}{rrrrrr}
\toprule
Phenotype & True $\pi_0$ & Bias & SE Ratio & Standard 95\% Cov. & Centered Cov. (Test A) \\
\midrule
1 & 0.2706 & $-0.0006$ & 1.04 & 94.8\% & 94.9\% \\
2 & 0.2579 & $-0.0042$ & 0.80 & 94.4\% & 96.0\% \\
3 & 0.2418 & $\hphantom{-}0.0012$ & 0.88 & 96.0\% & 95.5\% \\
4 & 0.2297 & $+0.0035$ & 1.12 & 96.9\% & 97.5\% \\
\bottomrule
\end{tabular}%
}
\caption{Validation results under the revised phenotype-probability clustering design ($c=0.3$), jackknife-based Godambe correction with full-chain retention, \emph{with} the design-degrees-of-freedom adjustment ($c_{\text{corr}} = M/(M-H)$ scaling of $\hat\Sigma_{\text{sandwich}}$ and a $t_{M-H}$ credible-interval critical value; $n=1{,}000$ Monte Carlo replicates). Mean $K_{\text{raw}} = 6.0$.}
\label{tab:cluster-probs-corrected}
\end{table}

\paragraph{Data-driven merging when $K_0$ is unknown.} Every result above relies on a merge-to-known-truth step (the fitted centroids nearest to the known $\text{MU}^0$ means) that is only available in simulation. On HCHS/SOL, $K_0$ is unknown, so this step must be replaced by an unsupervised rule operating on the MCMC output alone. We do so with a posterior-similarity-matrix point estimate of the unit-level partition: for each replicate, the co-clustering frequency across the (already label-switching-corrected, fully-retained) post-burn-in draws is summarized into an $n\times n$ posterior similarity matrix, and a point-estimate partition is obtained by maximizing the posterior expected adjusted Rand index of \citet{fritschickstadt2009}, an improved-Dahl criterion for exactly this problem, via the \texttt{mcclust::maxpear} implementation. The resulting estimated groups are used to build the marginal-likelihood density $a_{ik}$ and the merged prevalence functional exactly as in the truth-based pipeline above, but with the raw components mapped to \emph{estimated} groups rather than to known phenotypes; a centroid-distance alignment of the estimated groups to the $K_0$ known phenotypes is retained only to \emph{score} the method against ground truth in this validation study (analogous to computing an adjusted Rand index against known labels in any clustering simulation), never used by the merge itself.

Applied naively, this data-driven merge recovers the correct number of phenotypes, $\hat K=K_0=4$, in only $28.8\%$ of replicates, with a heavy right tail of over-segmentation (up to $17$ estimated groups in the worst replicate). We traced this to small, spurious satellite groups and pairs of nearby sub-clusters within a single true phenotype that nonetheless co-cluster distinctly often enough to register in the posterior similarity matrix. We suppress both failure modes with a second unsupervised post-processing step, applied to the estimated partition before any truth-based alignment: iteratively, either (i) the smallest estimated group is merged into its nearest neighbor by centroid distance if its total survey weight falls below $2\%$ of the total, or (ii) the two closest group centroids are merged if their Euclidean distance falls below $1.5$, repeating until neither condition triggers. This threshold recovers $\hat K = K_0 = 4$ in $98.2\%$ of replicates, with the remainder never falling below $4$ (i.e., no risk, in this design, of merging two genuinely distinct phenotypes together, since the true phenotype centroids are separated by at least $4.12$, well above the $1.5$ separation threshold).

Table~\ref{tab:datadriven-merge} reports validation metrics for all three data-driven variants against the truth-based ceiling (Table~\ref{tab:cluster-probs-corrected}). The size/separation threshold fixes $K_0$ recovery at essentially no cost to bias or coverage, and adding the design-degrees-of-freedom correction on top brings the fully unsupervised pipeline to $92.5$--$96.4\%$ coverage, within $1$--$2$ percentage points of the truth-based ceiling ($94.8$--$96.9\%$) on every phenotype, and with smaller bias on three of four phenotypes. We regard this as evidence that the survey-weighted DPMM pipeline, including its uncertainty correction, remains valid when $K_0$ is not known in advance, provided the raw-component merge is followed by this consolidation step. One caveat: the threshold values here ($2\%$ minimum weight, $1.5$ minimum separation) were chosen with reference to this design's true separation scale as a sanity check on the method, not tuned against ground-truth coverage; applied to HCHS/SOL, these thresholds would need to be set from domain knowledge of clinically meaningful biomarker separation or a data-driven heuristic (e.g., inspecting the merge sequence for a natural gap), rather than assumed to transfer directly from this simulation.

\begin{table}[h]
\centering
\small
\resizebox{\linewidth}{!}{%
\begin{tabular}{rrrrrrrr}
\toprule
& \multicolumn{2}{c}{No threshold} & \multicolumn{2}{c}{+ Threshold} & \multicolumn{2}{c}{+ Threshold + design-df} & \\
Phenotype & Bias & Cov. & Bias & Cov. & Bias & Cov. & True $\pi_0$ \\
\midrule
1 & $\hphantom{-}0.0009$ & 90.4\% & $\hphantom{-}0.0007$ & 89.6\% & $\hphantom{-}0.0006$ & 92.5\% & 0.2706 \\
2 & $-0.0025$ & 92.4\% & $-0.0024$ & 93.0\% & $-0.0015$ & 94.7\% & 0.2579 \\
3 & $\hphantom{-}0.0023$ & 93.0\% & $\hphantom{-}0.0017$ & 93.4\% & $\hphantom{-}0.0016$ & 96.4\% & 0.2418 \\
4 & $-0.0008$ & 93.5\% & $-0.0001$ & 93.5\% & $-0.0007$ & 95.7\% & 0.2297 \\
\bottomrule
\end{tabular}%
}
\caption{Validation results for the fully unsupervised, data-driven merge (no ground truth used in estimation; $c=0.3$, jackknife-based Godambe correction with full-chain retention, $n=1{,}000$ Monte Carlo replicates per column). ``Cov.'' is standard 95\% coverage. $K_0$ recovery rate: $28.8\%$ (no threshold), $98.2\%$ (+ threshold), $98.2\%$ (+ threshold + design-df, same partitions as the threshold column).}
\label{tab:datadriven-merge}
\end{table}

\bibliographystyle{unsrtnat}
\bibliography{references}

\end{document}
